\subsection{Japanese rings}

DEFINITION 1. --- Let A be an integral Noetherian ring. A is said to be Japanese if the integral closure of A in every finite extension of its field of fractions is a finite A-algebra.

Remarks. --- 1) It is equivalent to say that A satisfies the following condition: every integral domain B which is an A-algebra, integral over A, and contained in a finitely generated extension of the field of fractions K of A, is a finite A-algebra. Indeed, the field of fractions L of B is an algebraic extension of K, hence is of finite degree over K (A, V, p. 112, cor. 1 of prop. 17). The A-algebra B is contained in the integral closure of A in L, and is therefore finite if the latter is finite.

\begin{enumerate}
\def\labelenumi{\arabic{enumi})}
\setcounter{enumi}{1}
\tightlist
\item
  Let A be a Noetherian integral Japanese ring and S a multiplicative subset of A not containing 0. The ring of fractions \(\mathbf{S}^{-1}\mathbf{A}\) is Japanese. Indeed, let L be a finite extension of the field of fractions of A, and let B be the integral closure of A in L; then the integral closure of \(\mathbf{S}^{-1}\mathbf{A}\) in L is \(\mathbf{S}^{-1}\mathbf{B}\) (V, § 1, no. 5, prop. 16), hence is a finite \(\mathbf{S}^{-1}\mathbf{A}\)-algebra.
\end{enumerate}

Example. --- Every integral algebra of finite type over a field is a Japanese ring (V, § 3, no. 2, th. 2).

PROPOSITION 1. --- Let A be a Noetherian integral domain, and K its field of fractions. Suppose that for every finite radical extension L of K, the integral closure of A in L is a finite A-algebra. Then the ring A is Japanese.

Let E be a finite extension of K. Let N be a finite quasi-Galois extension of K containing E (A, V, p. 54, cor. 1), and let L be the field of invariants of the group of K-automorphisms of N. Then (A, V, p. 73, prop. 13), L is a radical extension of K and N is a separable extension of L. The integral closure B of A in L is therefore, by hypothesis, a finite A-algebra; the integral closure C of B in N is a finite B-algebra (V, § 1, no. 6, cor. 1 to prop. 18), hence a finite A-algebra. The integral closure of A in E is contained in C, and is therefore a finite A-algebra since A is Noetherian.

COROLLARY. --- Suppose the field K is perfect (for example, of characteristic 0). Then A is Japanese if and only if its integral closure is a finite A-algebra.

PROPOSITION 2. --- Let B be a Noetherian integral domain and A a Noetherian subring of B, such that B is a finite A-algebra. For A to be Japanese, it is necessary and sufficient that B be Japanese.

Let K (resp. L) denote the field of fractions of A (resp. B). Suppose first that A is Japanese, and let M be a finite extension of L. Let C denote the integral closure of B in M. By V, § 1, no. 1, prop. 6, C is the integral closure of A in M, hence is a finite A-algebra since M is a finite extension of K and A is Japanese. A fortiori, C is a finite B-algebra. This proves that B is Japanese.

Conversely, suppose B is Japanese and let N be a finite extension of K. Let D denote the integral closure of A in N. Let E be the compositum of L and N over K; since B is Japanese, the integral closure D' of B in E is a finite B-algebra, hence a finite A-algebra; the A-module D, which is a submodule of D', is therefore finitely generated, which implies that A is Japanese.

